\documentclass[12pt]{article}
\usepackage[utf8]{inputenc}
\usepackage[T1]{fontenc}
\usepackage{hyperref}
\title{Yakkha Language Resources for K-12 Curriculum Development}
\author{Compiled by Claude Sonnet 5.5 (claude-sonnet-5-5) for Andr\'as Kornai}
\date{}
\begin{document}
\maketitle

\section{Introduction}
Yakkha (ISO 639-3 ybh; Glottocode yakk1236) is a Sino-Tibetan (Kiranti) language of eastern Nepal and India. Estimates are about 10,000 speakers, and one source cites about 14,000. One source notes that the language is not being passed on to the next generation. The searches found strong descriptive work, but no school materials.

\section{Primary Educational Resources}
No primer, textbook or school programme for Yakkha turned up in the searches.

\section{Linguistic and Reference Materials}
\begin{itemize}
\item Diana Schackow, \emph{A Grammar of Yakkha} (Language Science Press, 2015, Studies in Diversity Linguistics 7). It is open access:
\url{https://langsci-press.org/catalog/book/66}
\item Open Textbook Library listing of the same grammar: \url{https://open.umn.edu/opentextbooks/formats/2010}
\item Source repository of the grammar on GitHub: \url{https://www.github.com/langsci/66}
\item OAPEN record of the grammar: \url{https://0-library-oapen-org.catalogue.libraries.london.ac.uk/handle/20.500.12657/32844}
\item Max Planck Institute for Evolutionary Anthropology project page for the Yakkha grammar: \url{https://www.eva.mpg.de/linguistics/past-research-resources/documentation-and-description/a-grammar-of-yakkha/}
\end{itemize}

\section{Religious and Community Resources}
No religious or community-organisation resources turned up in the searches.

\section{Digital and Audio Resources}
\begin{itemize}
\item ELDP archive collection ``Documentation and grammatical description of Yakkha, Nepal''. One source describes it as including an annotated audio corpus and a trilingual dictionary; I did not confirm what is publicly available: \url{https://eldp.access.preservica.com/dk0293/}
\end{itemize}

\section{Curriculum Development Context}
A comprehensive open-access grammar and an archived corpus exist, which is a solid base for building native-language materials. No orthography in school use or teaching programme was found. Any curriculum work would need partnership with the Yakkha community.

\end{document}